
\documentclass{article}
\usepackage{colortbl}
\usepackage{makecell}
\usepackage{multirow}
\usepackage{supertabular}

\begin{document}

\newcounter{utterance}

\centering \large Interaction Transcript for game `cladder', experiment `full\_v1.5\_default', episode 8201 with selfrefine{-}llama{-}3.1{-}8b{-}instruct.
\vspace{24pt}

{ \footnotesize  \setcounter{utterance}{1}
\setlength{\tabcolsep}{0pt}
\begin{supertabular}{c@{$\;$}|p{.15\linewidth}@{}p{.15\linewidth}p{.15\linewidth}p{.15\linewidth}p{.15\linewidth}p{.15\linewidth}}
    \# & \multicolumn{2}{c}{Player} && \multicolumn{2}{c}{Game Master} \\
    \hline

    \theutterance \stepcounter{utterance}  
    & & & \multicolumn{4}{p{0.6\linewidth}}{
        \cellcolor[rgb]{0.9,0.9,0.9}{
            \makecell[{{p{\linewidth}}}]{
                \texttt{\tiny{[P1$\langle$GM]}}
                \texttt{You are an expert in causal inference. The following question is not a typical commonsense query, but rather a meticulously designed question created by a professor specializing in causal inference, intended to assess the students' mastery of the course content.} \\
\texttt{Imagine a self{-}contained, hypothetical world with only the following conditions, and without any unmentioned factors or causal relationships: Tijv has a direct effect on tijw and xevo. Xevo has a direct effect on gyzp. Tijw has a direct effect on gyzp. For those who are not xevo and are not tijw, the probability of gyzp is 76\%. For those who are not xevo and are tijw, the probability of gyzp is 76\%. For those who are xevo and are not tijw, the probability of gyzp is 12\%. For those who are xevo and are tijw, the probability of gyzp is 14\%. For those who are not xevo, the probability of tijw is 28\%. For those who are xevo, the probability of tijw is 14\%. For those who are xevo, would it be more likely to see gyzp if the individual was not xevo?} \\
\texttt{Start your answer with 'Yes' or 'No', followed by additional reasoning or evidence to support your explanation.} \\
            }
        }
    }
     \\ \\

    \theutterance \stepcounter{utterance}  
    & \multicolumn{4}{p{0.6\linewidth}}{
        \cellcolor[rgb]{0.561,0.737,0.561}{
            \makecell[{{p{\linewidth}}}]{
                \texttt{\tiny{[P1$\rangle$GM]}}
                \texttt{No.} \\
\\ 
\texttt{To determine whether it would be more likely to see gyzp if the individual was not xevo, we need to compare the probability of gyzp if the individual is not xevo (Y(0)) to the probability of gyzp if the individual is xevo and not tijw (Y(1) | W(1) = 0). } \\
\\ 
\texttt{We have already calculated P(Y(0) = 1) = 0.76. We also know that P(Y(1) = 1 | W(1) = 0) = 0.12. } \\
\\ 
\texttt{To compare these probabilities, we need to calculate the probability of gyzp if the individual is xevo and not tijw, which is P(Y(1) | W(1) = 0). However, we are given P(Y(1) = 1 | W(1) = 0) = 0.12, which is the probability of gyzp if xevo and not tijw, but not the probability of being xevo and not tijw. } \\
\\ 
\texttt{Since we are given that P(Y(1) = 1 | W(1) = 1) = 0.14 and P(Y(1) = 1 | W(1) = 0) = 0.12, we can see that the probability of gyzp if xevo and tijw is higher than the probability of gyzp if xevo and not tijw. } \\
\\ 
\texttt{However, we are interested in the probability of gyzp if the individual is xevo and not tijw, which is P(Y(1) | W(1) = 0). To calculate this, we need to know the probability of being xevo and not tijw, which is P(W(1) = 0 | Xevo = 1). } \\
\\ 
\texttt{Unfortunately, we are not given this probability. However, we can use the law of total probability to write:} \\
\\ 
\texttt{P(W(1) = 0 | Xevo = 1) = P(W(1) = 0 | Xevo = 1, Tijw = 0)P(Tijw = 0 | Xevo = 1) + P(W(1) = 0 | Xevo = 1, Tijw = 1)P(Tijw = 1 | Xevo = 1)} \\
\\ 
\texttt{We are given that P(Tijw = 0 | Xevo = 1) = 1 {-} P(Tijw = 1 | Xevo = 1) = 1 {-} 0.14 = 0.86, and P(Tijw = 1 | Xevo = 1) = 0.14. } \\
\\ 
\texttt{We are also given that P(W(1) = 0 | X} \\
            }
        }
    }
    & & \\ \\

    \theutterance \stepcounter{utterance}  
    & & & \multicolumn{2}{p{0.3\linewidth}}{
        \cellcolor[rgb]{0.9,0.9,0.9}{
            \makecell[{{p{\linewidth}}}]{
                \texttt{\tiny{[GM$|$GM]}}
                \texttt{No} \\
            }
        }
    }
    & & \\ \\

    \theutterance \stepcounter{utterance}  
    & & & \multicolumn{2}{p{0.3\linewidth}}{
        \cellcolor[rgb]{0.9,0.9,0.9}{
            \makecell[{{p{\linewidth}}}]{
                \texttt{\tiny{[GM$|$GM]}}
                \texttt{yes} \\
            }
        }
    }
    & & \\ \\

    \theutterance \stepcounter{utterance}  
    & & & \multicolumn{2}{p{0.3\linewidth}}{
        \cellcolor[rgb]{0.9,0.9,0.9}{
            \makecell[{{p{\linewidth}}}]{
                \texttt{\tiny{[GM$|$GM]}}
                \texttt{game\_result = LOSE} \\
            }
        }
    }
    & & \\ \\

\end{supertabular}
}

\end{document}
